\chapter{A shared service for economic analysis}
\label{r3:product}

A banking study must connect a business choice with a financial consequence
through a defensible account of customer behaviour and observed data. A commercial question names a feasible action
and a financial consequence. Its data are event records with changing
identities and reporting delays. An economist supplies restrictions on
assignment or shocks; a learner estimates conditional distributions; a
validator needs a sampling argument; a decision maker needs the value of
a feasible policy. A shared product is useful when it preserves the
meaning of the question across those transformations.

The proposed product is a library and a standalone analytical workspace
built around the same definitions of populations, interventions and outcomes. Its first release should expose
a bounded collection of well-defined operations. More general discovery,
inverse problems and latent control remain important research capabilities,
but their outputs must retain the assumptions and unresolved alternatives
developed in the preceding chapters. The construction below derives this
boundary; it is a specification, not a claim that the entire product
already exists.

\section{Define the study before choosing an estimator}
\label{r3:request-type}

Write a request as
\begin{equation}
 \mathcal Q=(\mathcal M,\mathcal K,g,\nu,\theta,\mathcal E,\mathcal F),
 \label{r3:request}
\end{equation}
where $\mathcal M$ is a class of causal models, $\mathcal K$ the observation
operator, $g$ the intervention, $\nu$ the target population, $\theta$
the requested functional, $\mathcal E$ the sampling experiment and
$\mathcal F$ the information available to the decision rule. For a model
$M\in\mathcal M$, the observed law is $P=\mathcal K(M)$ and the causal
answer is $\theta(M;g,\nu)$. The observation operator includes selection,
measurement, clocks, aggregation and record linkage when relevant.
The sampling experiment specifies independent households, randomized
branches, a dependent calendar series, or another explicit design.

For deposit retention, the request might ask for the difference in
six-month discounted net funding contribution under two invitation
rules for households eligible on a specified date. It must distinguish
invitation from receipt, and baseline eligible households from eventual
activators. If the action is instead an increase in offered interest,
both customer balances and the price paid on existing balances change.
These are different interventions and financial functionals even when
the same table is used.

The target needs a measurable definition before computation. For example,
with monthly contribution $R_t$ in a declared currency and discount
factor $d_t$, use
\[
 \theta(M;g_1,g_0,\nu)
 =\E^\nu_M\!\left[\sum_{t=0}^{T-1}d_tR_t^{g_1}\right]
  -\E^\nu_M\!\left[\sum_{t=0}^{T-1}d_tR_t^{g_0}\right].
\]
Here $R_t$ includes the agreed funding, servicing, incentive and expected
loss components without counting the same balance benefit twice.
An average balance effect is an input to this account, not already its
answer. Currency conversion, discounting, exposure weights and eligibility
are part of the target rather than invisible formatting options.

Several request fields can be inferred syntactically from data, but the
scientific restrictions cannot. Software can detect a duplicate decision
or a covariate first recorded after action. It cannot infer that a pricing
instrument has no direct effect or that all private underwriting
information has been measured. The workspace should let an analyst draw
and explain the model, compare plausible alternatives, and attach empirical
diagnostics to particular restrictions. A single numerical validity score
would erase the distinction between those obligations.

\section{Identification returns a mathematical object or an explicit limit}
\label{r3:identification-interface}

Define the identified set
\[
 \Theta_{\mathcal Q}(P)=
 \{\theta(M;g,\nu):M\in\mathcal M,\ \mathcal K(M)=P\}.
\]
An identification routine receives the model and observation declarations;
it does not receive permission to add whichever assumptions make a preferred
estimator available. A singleton set can be represented by an observed-law
functional $\Psi(P)$. A non-singleton set may admit bounds, a union of
functionals, or only a nonidentification witness. The useful result types
are therefore distinct.

\begin{description}[style=nextline,leftmargin=0pt,labelwidth=0pt,labelsep=0pt]
\item[Identified expression.] An expression for $\Psi$ together with the
assumptions and transformations that derive it. Numerical evaluation may
still be unsupported.
\item[Identified set or outer bound.] A set with its model class, support,
restrictions and sharpness status. A conservative outer bound is useful,
but must not be labelled sharp without a compatible-model construction.
\item[Nonidentification witness.] Two allowed models with the same observed
law and different targets, or a valid source-specific failure certificate
such as a hedge in the applicable ID model.
\item[Unsupported query.] The registered algorithm does not handle the
graph class, intervention, observation process or requested target.
This makes no claim that the target is unidentified.
\item[Incompatible specification.] The declarations are logically
inconsistent or violate a deterministic support/schema requirement.
A statistical model-check rejection should separately carry its test
and error level.
\end{description}

These types follow from Chapter~\ref{r3:graphs} and the ID/IDC construction
of \citet{id2006}. Failure to find a back-door set is not a
nonidentification certificate. Nor is failure of an optimizer. For
example, the front-door expression is identified under its restrictions,
but a backend implementing only covariate adjustment should return
an identified expression with no available estimator. Conversely, the
two invitation models of Section~\ref{r3:late} agree on all observed
records and disagree on the population ATE. Increasing an estimator's
iteration budget cannot supply that missing information.

For an initial implementation, observed-law expressions can use a small
typed language: conditional probability kernels, products, integrals or
finite sums, substitutions, and ratios on declared positive support.
A kernel records its random arguments, conditioning arguments, units,
dominating measure and observation regime. The type of
$p(y\mid a,x)$ differs from $p^g(y\mid x)$; a density in currency units
differs from a probability mass. Marginalization removes a random argument;
it must not accidentally integrate a fixed intervention value.

Each allowed transformation should produce a trace. In ordinary adjustment,
that trace names the adjusted variables, explains why their predecision
status and separation condition apply, and yields
$\int p(y\mid a,x)p(x)\,dx$. In recursive ID, it records ancestors,
districts, extracted kernels and support restrictions. The trace certifies
the consequence of the supplied assumptions. It cannot certify that
the graph supplied by a bank economist is the true graph.

\section{Adapters must preserve targets and statistical experiments}
\label{r3:adapter-inference}

The target-preservation identity
$\Psi(P)=\Psi_B(T_\#P)$ in Section~\ref{r3:software-types} is necessary
for converting data through $T$ to backend $B$. It also constrains
uncertainty. Suppose, for this paragraph, that observations $O_i$ are
independent and identically distributed, $T$ is fixed and applied
separately to each unit, and $Z=T(O)$. Let $P_\varepsilon$ be a regular
submodel with score $S(O)$. Differentiating the pushforward density gives
the reduced score
\[
 S_Z(Z)=\E[S(O)\mid Z].
\]
If $\phi_B(Z)$ is a gradient for the reduced functional, then
\begin{align}
 \left.\frac{d}{d\varepsilon}\Psi(P_\varepsilon)\right|_0
 &=\E[\phi_B(Z)S_Z(Z)]\nonumber\\
 &=\E[\phi_B(T(O))S(O)].
 \label{r3:gradient-pullback}
\end{align}
Thus the pulled-back score is a gradient for the original target.
In an unrestricted original model it is the canonical gradient when
the stated differentiability conditions hold; in a restricted model
projection onto the model's tangent space may improve efficiency.
This is a project consequence of the pathwise argument in
Chapter~\ref{r3:learning}, not a theorem for arbitrary preprocessing.

The qualifications explain common failures. A learned encoder is random
and may need independent training and a remainder analysis. Demeaning
over a panel is not a fixed transformation of each independently sampled
row. Expanding one household into transactions changes the sampling
units. Discarding a component of history may destroy target preservation.
An adapter must therefore declare both its mathematical transformation
and its mapping of sampling units, weights and information sets.

A backend estimator should return sufficient information for the
specified inference route. For a regular scalar moment this includes
unit-level scores, the Jacobian, target normalization, fold assignments
and nuisance estimates or their immutable references. For an identified
set it includes bound objectives and constraint certificates. For a
weak-IV inversion it includes the accepted set, including disconnected
or unbounded components. For a likelihood model it includes the exact
factorization, numerical approximation and parameter restrictions.
Forcing all of these outputs into an estimate/standard-error pair would
discard information that changes the bank decision.

\label{r3:protocol}
The companion scientific protocol gives this separation a concrete design
form. A request binds the model, observation, target and sampling
specifications by content hash. A result separately records computational
completion, identification and inferential availability. A scalar
confidence set is an ordered union of intervals, including empty,
disconnected and unbounded sets; posterior credible sets have a different
type. Derivations, bounds and failure witnesses have distinct artifact
roles. The supplied schema fixtures check these representation invariants.
They do not execute an identification algorithm or establish coverage.
The following versioned extension supplies explicit vector, curve and graph-set
types for the bounded constructions developed in this monograph. A general
graph probability distribution remains a different object.

\section{A bounded mathematical language and a concrete lowering rule}
\label{s4:scientific-language}
The specification must be stronger than a collection of file hashes. A
variable declaration contains its name, support, physical unit, causal role
and clock. A population declaration identifies the economic unit and the
role of its weights. An observation declaration identifies the immutable
view, revision relation, availability convention and sampling frequency.
The model specifies a structural class and an assumption set. A target
specifies the action, outcomes, fixed horizons or policies, normalization
and population. These declarations together determine a statistical
experiment; no single field can substitute for the others.

For finite mechanisms an executable expression grammar is particularly
transparent. A mass kernel has ordered parent variables $U$, ordered child
variables $V$, declared finite supports and an array of shape
$|\mathcal X_U|\times|\mathcal X_V|$. Every parent slice must be
nonnegative and sum to one over the child axes. A generative product
$\prod_jK_j(x_j\mid x_{\operatorname{pa}(j)})$ requires one mechanism
per child and a topological order. Marginalization sums named random axes;
substitution fixes an axis without summing it. Replacing an action mechanism
by $g(a\mid h)$ requires that its information variables precede the action.
This operation evaluates the supplied causal factorization; it does not
establish that an arbitrary observed kernel can be intervened upon.

The local finite-mass interpreter implements product, marginalization and
mechanism replacement, checking support, normalization and information order.
For a binary confounded example it recovers the interventional mean $0.5$
and rejects a policy using a future outcome. A general expression language
would additionally need density measures, positive-denominator restrictions
and numerical integration error. The initial lowering rules do not pretend
that this finite interpreter implements continuous ID or general latent
identification. Instead they recognize the complete constructions already
derived and reject mismatched target/sampling combinations.

\begin{longtable}{P{.13\textwidth}P{.38\textwidth}P{.42\textwidth}}
\toprule Regime & Scientific input & Lowering and inferential output\\\midrule\endhead
Stationary response & Fixed horizon, stable Gaussian linear law, external
instrument and nonzero impact normalization. & Joint direct moments,
Bartlett covariance and residual-test interval unions; a proxy VAR is a
challenger only with its additional dynamics restriction.\\
Integrated response & Stationary changes, same shock restriction, cumulative
level target. & Difference, fit common-origin moments, transform their full
covariance, then invert. No stationary-level shortcut.\\
Observed history & Independent whole trajectories, sequential identification,
fixed feasible policies and the finite observed-state fitting family. &
Cross-fitted augmentation or backward targeting; within-customer policy
contrast scores and trajectory variance.\\
Controlled latent & Randomized logger, binary informative emission model,
identified transitions/rewards and regular likelihood. & Scaled EM plus
observed-score refinement; fixed-policy value and information/refit uncertainty.
Boundary or failed fits retain an unavailable inference status.\\
Gaussian graph class & Three observed variables, iid Gaussian sample,
causal sufficiency and optional background order. & Joint covariance region,
retained DAGs and outer effect union. No probability attached to a DAG.\\
Temporal graph class & Stable VAR(1), known lag bound, independent innovations,
no contemporaneous causal edges. & Coefficient box, unresolved zero/nonzero
edges and its fixed-horizon response image.\\\bottomrule
\end{longtable}

These are conditional method choices, not six interchangeable backends.
The standalone reference runner checks the request and data identities,
dispatches these bounded numerical recipes and validates their typed result.
Its two servicing families retain their stated reward, state and policy
semantics; a new family requires a new mathematical and numerical qualification.
The separated-block sparse IV construction is available as a reproducible
study recipe with declared learner settings. General controlled continuous
states, unrestricted temporal discovery and policy optimization remain
separate extension requests.

\subsection{Results that can carry the actual uncertainty}
Protocol version 0.4 supplements the historical scalar protocol. A scalar
is a one-coordinate result; vectors and response curves have named ordered
axes and one physical unit per flattened coordinate. A response curve orders
outcome before horizon. Its covariance has that same coordinate order and
already includes the sample-size divisor. A confidence object carries its
method, level, coverage target and an interval union for each coordinate.
For a simultaneous rectangular set, the Cartesian product of these unions
is the confidence region. Empty, disconnected and unbounded coordinate sets
remain representable; an unbounded endpoint is a typed null, not a JSON
infinity or a large arbitrary number. Pointwise coverage cannot be relabeled
simultaneous by an adapter.

A graph result retains ordered nodes and its graph class. A DAG set stores
directed binary adjacency matrices. CPDAG/PAG endpoint matrices use zero
for absence, one for a tail, two for an arrowhead and three for a circle;
the latter is permitted only for a PAG. Basic endpoint validation is not
a complete PAG-recognition algorithm. The temporal reference instead stores
lower and upper coefficient matrices with edge marks ``present'' or
``zero or nonzero''. Its graph class is not a posterior distribution over
graphs. The effect union is explicitly an outer set, and uncertainty about
the causal-sufficiency assumption is not encoded as its confidence level.

Computational completion, conditional identification and inference have
separate status fields. Numerical error bounds, when available, use the
same output units. A null numerical bound means no certified bound was
provided; it must not be interpreted as zero error. Solver diagnostics and
all requested failed repetitions remain attached evidence. Hashes bind
those arrays and declarations to a computation, while scientific validation
determines whether their asserted premises are plausible.

\subsection{Allowed conversions and explicit rejection cases}
For a fixed known diagonal scale matrix $S$, the converted effect and
covariance are $S\widehat\theta$ and $S\widehat V S'$. Interval endpoints
reverse when a scale is negative; the order of disconnected components
also reverses. A rate denominator changing from percentage points to basis
points divides a normalized response by 100. A fixed currency multiplier
can be applied in the same way. A time-varying exchange rate generally
changes the reward functional and requires rebuilding it; a monthly stock
cannot be converted into daily flows by relabeling an axis.

Population weights, inverse behavior probabilities, inverse observation
probabilities and bootstrap frequencies have different mathematical roles.
The request validator rejects their substitution. It also rejects an iid
row experiment for the dependent response route, a split inside a purportedly
independent trajectory, inclusive availability when the request requires
strict predecision information, and a learned policy supplied to fixed-policy
inference. A valid clock conversion can preserve timestamp/sequence order
exactly; changing that order requires a new observation analysis. These
checks enforce declared meaning, not the truth of exchangeability or
instrument exclusion.

\section{Concrete sibling boundaries and the standalone study}
\label{s4:sibling-mappings}
Read-only inspection of the actual sibling APIs changes the integration
design. FinMarketing's \texttt{CausalData.from\_record} requires pretreatment
features, mature outcomes, binary actions and independent groups. Its
\texttt{FeatureSnapshot} orders availability, cutoff and decision by
timestamp/sequence pairs. The local mapping preserves those pairs, horizon
and dependence unit. It does not coerce a dynamic curve into FinMarketing's
scalar \texttt{Estimate}; that requires a richer result consumer. A causal
response to a macro shock also cannot inherit its independent-group standard
error merely because the source rows are account records.

dataml's \texttt{LineageGraph.append} accepts run, dataset, partition and
metadata, while its correction schema requires a resolvable assumption and
a numerical confidence field. Observation candidates therefore map to
lineage metadata with their assumptions and alternatives. An uncalibrated
hint score does not populate correction confidence, and ambiguous candidates
do not generate an automatic correction. This preserves dataml's distinction
between an original source and an assumed overlay.

BayesFilter's \texttt{SSMStaticShape}, \texttt{SSMDataSignature} and
\texttt{BayesianSSMProblem} can record dimensions, data/mask identities and
parameter coordinates. For the binary controlled family, the local mapping
adds action identity, randomized-logger semantics, reward, finite state
support and the full observation order to its model manifest. A generic
Bayesian SSM descriptor alone does not supply those causal meanings or a
controlled categorical filter implementation. Its numerical target and
derivative still need to refer to the same likelihood, as required by
Chapter~\ref{r3:controlled-identification}.

MacroFinance's \texttt{DirectedPairMeasurement} and
\texttt{ObservationContract} distinguish directed quote pairs, maturity,
observed rows and derived diagnostics. Its convention is $x_m=-\chi$;
reversing a pair reverses the quoted wedge sign. The local mapping preserves
these identities and requires a separate causal role for the measured
variable. A triangle diagnostic is not automatically an identified structural
cycle intervention. The inspected MacroFinance BayesFilter facade also
references a different import root from the workspace checkout, so inspection
alone cannot establish runtime compatibility between them.

The same scientific request/result files support an independent study:
choose the economic target and observation plan, inspect the identifying
assumptions, select a supported regime, execute its estimator and mandatory
comparison, and publish the result with its uncertainty and source trace.
Local examples exercise the mappings and numerical operations here. Sibling
checkouts remain unchanged; accepted adoption, shared release management and
institutional data access remain separate implementation work.


\section{A bounded catalogue of scientific operations}
\label{r3:operation-catalogue}

The first supported catalogue should be selected by mathematical closure
and banking demand. Table~\ref{r3:operations} defines seven operation
families. Support means a particular target, observation representation,
estimator and inference construction have been qualified together;
it does not mean every method within that family is implemented.

\begin{longtable}{P{.16\textwidth}P{.38\textwidth}P{.39\textwidth}}
\caption{Proposed operation families and their required scientific outputs}
\label{r3:operations}\\
\toprule
Operation & Input and identification route & Output and essential rejection\\
\midrule\endfirsthead
\toprule Operation & Input and identification route & Output and essential rejection\\
\midrule\endhead
Adjusted effect &
One decision, baseline covariates, explicit exposure/population weights;
consistency, conditional exchangeability and support. &
Regime mean or ATE functional, cross-fitted score and declared cluster
inference. Reject a descendant adjustment or an unsupported action.
CATE surfaces need their own local inference route.\\
Invitation/IV &
Separate assignment/receipt; randomized or justified conditional instrument;
linear moments, monotone compliance or a stated inverse model. &
ITT, homogeneous coefficient or weighted complier target. Preserve weak-IV
test sets; do not relabel a LATE or regularized response as an ATE.\\
Adoption and panel contrasts &
Entity/calendar records; cohort, comparison group, anticipation and
parallel-trend specification, or a distinct structural moment model. &
Group--time ATT and declared aggregation, or the named panel coefficient.
Reject conversion of reversible actions into first-adoption dates.\\
Observed-history policy &
Finite horizon, ordered histories, feasible actions, intervention density,
event/censoring definitions and sequential identification restrictions. &
Policy value/contrast with its recursive estimator and trajectory score.
Differentiate propensity-dependent policies; retain horizon/support limits.\\
Dynamic shock response &
Observation frequency, moment restrictions, shock instrument and scale,
dependence regime and a finite set of horizons. &
Normalized response vector with stacked inference or weak-normalization
sets. Reject an unidentified historical-shock interpretation.\\
Controlled latent analysis &
Initial, transition, emission, reward and logging laws; policy information
and a target-specific identification argument. &
Model-conditional fit/filter, identified value or value ambiguity.
Planning requires the interventional law, not only a fitted likelihood.\\
Discovery-to-query &
Clock, graph class, background constraints, CI/functional assumptions and
the observation operator. &
Endpoint-preserving equivalence graph, then a common functional, effect set
or unresolved query. Retain uncertainty about discovery itself.\\
\bottomrule
\end{longtable}

For adjusted and longitudinal means, the product can initially support
the scores derived in Chapters~\ref{r3:learning} and \ref{r3:longitudinal},
with finite categorical actions and fixed horizons. This is a useful
construction boundary: it permits both simple randomized comparisons and
repeated collections policies while keeping the intervention density
explicit. A continuous deterministic intervention under a continuous
observed action law is not automatically absolutely continuous; density
ratios for point masses require a different identification/estimation
route. Modified policies and smooth stochastic shifts can supply one,
under the assumptions in Chapter~\ref{r3:policy-targets}.

Panel support should begin with specified cohort estimands and strong
linear-moment baselines, not an unrestricted auto-panel operation.
The equal-weight branch target, household target and balance-weighted
target require different aggregations. Temporal shock support should
likewise identify a normalized response under the stated instrument
exclusions, with the fixed-horizon inference of
Section~\ref{r3:series-inference}. Structural breaks and weak impact
normalizations should produce an explicit alternative analysis.

Controlled latent support can be valuable before a general latent causal
identification engine exists. The same workspace can compare a randomized
logging model, a private-logging model and an observed-history policy
value, displaying which conclusions require which assumptions. The
known-emission inversion in Section~\ref{r3:known-emission} supplies a
finite validation case; it is not a general theorem for learned emissions.
Proxy-value routes should be registered with their bridge and rank
conditions rather than offered through a generic hidden-state option.

\section{Estimation and uncertainty form one qualified composition}
\label{r3:inference-contract}

An inference contract has five parts: the target and its differentiability
class; an estimator expansion; the statistical sampling units; the
covariance or test construction; and the conditions controlling nuisance
and numerical error. For a regular vector target one common form is
\[
 \widehat\theta-\theta
 =-\widehat J^{-1}\frac1n\sum_i\psi_i(\theta,\eta_0)
   +R_n,\qquad \sqrt n R_n=o_p(1),
\]
where the meaning of $n$ and the scaling must match the experiment.
Independent trajectories, independent clusters with variable sizes,
crossed arrays and a long time series imply different normalizations
and limit arguments. A covariance option alone cannot establish them.

For one-way independent clusters, reduction of all score contributions
within each cluster precedes forming outer products. For crossed
entity/date dependence, inclusion--exclusion of cluster covariance
components requires the appropriate crossed-array assumptions; a large
entity count does not compensate for a fixed handful of macro dates.
Multiway cross-fitting must prevent training/evaluation overlap in the
relevant coordinates \citep{cluster2021}. A simple entity split is not
automatically a multiway construction.

For a time series, keep the time order, outcome window, separation gap
and score bandwidth as distinct fields. Overlapping outcomes can extend
dependence past disjoint origin blocks. When a stationary mixing
argument is invoked, its mixing and moment restrictions must be named.
A user selecting a HAC option has specified an estimator family, not proved
that a structural-break or unit-root score obeys the required CLT.
Simultaneous bands require a joint horizon construction.

Uncertainty for weakly identified parameters, nonsmooth bounds and
selected policies needs separate operations. Ratios with weak
denominators should use a qualified test inversion; graph ambiguity
requires a union or other defensible set operation. Learned-policy
evaluation must distinguish a frozen policy on an independent evaluation
sample from a data-adaptive optimal-value target. Bayesian posterior
intervals require their likelihood, prior and numerical diagnostics;
they cannot inherit frequentist coverage merely because they are called
credible.

The result should consequently carry several evidence dimensions:
identification conditional on stated restrictions, numerical convergence,
sampling inference, sensitivity to disputed restrictions and decision
adequacy. A converged point estimate can coexist with unsupported
uncertainty. A wide but valid confidence set can be adequate for a decision
if every contained model favors the same affordable action.

\section{Policy evaluation has a separate information boundary}
\label{r3:planner-interface}

A controlled-state adapter must expose the observational and interventional
laws separately. In the finite model, a pre-action belief is
$b_t(s)=P(S_t=s\mid H_t)$. Under a private logger, after seeing the
historical action it becomes proportional to
$b_t(s)q_\eta(A_t\mid s,H_t)$ before the transition, reward and next
signal update. Under an intervention using only the recorded history,
the action density $g_t(A_t\mid H_t)$ is constant across hidden states
and cancels in normalized filtering. Replacing the logger changes the
law, not simply one input tensor.

The planner receives a distribution over the information it is allowed
to use. A smoothed posterior $P(S_t\mid H_T)$, $T>t$, includes future
evidence and is not an online policy state. A retrospective smoother is
useful for fitting and diagnosis, while prospective evaluation must
reconstruct the predecision filter. The adapter must also specify whether
reward is observed before the next signal and whether reward carries
additional information about the state.

For a supplied causal model the finite-horizon value recursion is
\[
 V_t(b)=\max_{a\in\mathcal A_t(b)}
 \left\{r_t(b,a)+
 \gamma\int V_{t+1}(\tau_t(b,a,o',r))\,
                 p_t(do',dr\mid b,a)\right\}.
\]
The belief update $\tau_t$ and predictive observation/reward law must
use the same factorization. This is the information-state construction
of \citet{kaelbling1998} specialized to the declared timing.
For this display all decision-relevant recorded context is assumed to
be included in the belief state or its explicit arguments. Action
constraints can depend on recorded eligibility, customer consent
and a specified budget. An estimate of a hidden risk state does not
authorize using unavailable information in those constraints.

Model ambiguity creates another decision problem. If $\mathcal M_P$
contains several observationally compatible causal models, a robust
policy may maximize $\inf_{M\in\mathcal M_P}V_M(\pi)$. In general the
same $M$ must govern the whole trajectory. A stagewise minimization
that freely chooses a new adverse model at each date solves a larger,
rectangular uncertainty problem. Its result can be conservative and
need not be the sharp value bound for persistent hidden types. The
dynamic-bound analysis of Section~\ref{r3:dynamic-bounds} shows why
the uncertainty-set type must accompany the planning algorithm.

For a two-date fixed policy with no discount, suppose the allowed
reward vectors are $(0,1)$ and $(1,0)$, with one model applying at
both dates. Worst-case total value is one. Independently choosing
the worse component at each date instead gives zero, a vector absent
from the original model class. This exact example isolates the extra
freedom introduced by rectangularization; it is not a claim that a
bank's ambiguity set contains those two models.

Finally, solver error needs its own budget. The Bellman residual bound
in Section~\ref{r3:value-errors} controls numerical error for the
specified contraction. Parameter uncertainty, model ambiguity, reward
accounting and policy support remain separate. A planner can report
a precise solution to a scientifically weak model; the user needs
enough information to recognize that situation.

\section{Discovery should improve a decision-relevant model class}
\label{r3:discovery-product}

The discovery workspace begins with mechanisms and timing. It proposes
variables and lags from the available data, records business restrictions,
and returns a DAG, CPDAG, PAG or named functional-model object with its
endpoint semantics intact. Metadata lineage is displayed separately
from causal ancestry. An upstream database dependency does not establish
an economic intervention edge.

The next operation asks a specific causal query across compatible models.
For the chain example of Section~\ref{r3:graph-effects}, the observed
Gaussian covariance and CPDAG imply effects $\{0,1/4\}$ within the
declared linear causally sufficient class. The product should expose
both values, explain the orientations that produce them, and show
which credible timing information or intervention would distinguish
them. Returning the effect of one arbitrarily oriented representative
would turn an unresolved scientific question into a spurious estimate.

General PAG completion is expensive and may still leave latent
nonidentification within each completion. The first bounded route
should therefore use verified generalized-adjustment criteria on
supported graph classes \citep{perkovic2018}, small enumerated classes
for exact examples, and unresolved results beyond those capabilities.
Search limits, CI test assumptions, variable screening and lag windows
must be retained. A timeout is computationally unresolved, not a
certificate that no adjustment set exists.

A useful experiment recommendation should name the ambiguity it resolves.
For deposit pricing, a randomized encouragement may distinguish a price
response from latent liquidity selection. For collections, a new
predecision signal may improve a bridge's measurement rank. For treasury,
a higher-frequency observation may remove an aliasing ambiguity. Their
cost and decision value can be compared; edge-count improvement alone
is not the economic objective.

\section{Acceptance is a collection of scientific obligations}
\label{r3:product-acceptance}

Each registered operation needs a reviewable construction dossier. It
contains the exact target and restrictions, an identification derivation
or witness, an estimator/inference derivation, the qualified backend
version and transformation, and reproducible positive and negative
examples. Exact finite laws should be used when they expose semantic
errors that a large benchmark can conceal.

For adjustment, compare the estimator's population expression with
explicit intervention enumeration, and include a confounded pair of
models that observation cannot distinguish. For IV, hold the observed
law fixed while changing the unobserved outcomes of never-takers; ensure
that the requested ATE remains unresolved although LATE is returned.
For longitudinal analysis, verify the recursive score under feedback,
then change the policy's dependence on the propensity and require the
corresponding gradient change.

For controlled states, enumerate hidden paths to check the likelihood,
filter and sufficient counts with and without the logger; compare
interventional value against independent dynamic programming. For
discovery, round-trip all endpoint marks and propagate every supported
completion into its effect set. For data organization, delay and revise
events across the decision boundary and verify exactly which history
the policy receives. For distributed inference, repartition the same
scores across workers and preserve folds, cluster sums and accepted
uncertainty sets.

Repeated-sample experiments then assess bias, coverage or policy regret
under named laws, including failures of support, rank, stationarity
and model specification. Their purpose is to test the implemented
construction and characterize its operating range. They do not prove
that the bank's unobserved confounding assumption holds. Bank data
validation subsequently evaluates the actual information, process
and financial account.

This division makes a staged product possible without shrinking its
scientific ambition. A bank can use rigorous observed-history analyses
while researchers develop controlled latent and discovery-to-effect
routes in the same representations. Promotion of a method then requires
its missing mathematical and empirical obligations to be discharged,
rather than a new wrapper around an attractive estimator name.

\section{Executable composition across distinct scientific questions}
\label{s5:composition}
The common layer should earn its cost by preserving meaning across reused
operations. A small implemented composition begins with an
\texttt{AdjustmentQuery}: action, outcome, controls, population unit,
outcome units and sampling law. A supplied acyclic graph is checked using
the sufficient back-door criterion: reject action descendants among controls,
remove outgoing action arrows, take the ancestral moral graph and test
separation after removing the controls. A passing trace determines the
functional
\[
 \psi=\E_X[m_1(X)-m_0(X)]
\]
and the nuisance roles $(m_0,m_1,e)$. The graph is supplied scientific
knowledge; data do not certify its absent edges. Failure of this sufficient
criterion means this adjustment backend cannot identify the query, not that
no other identification strategy exists.

Whole-unit folds produce out-of-fold predictions. A separate score operation
forms
\[
 U=m_1(X)-m_0(X)+\frac{A\{Y-m_1(X)\}}{e(X)}
 -\frac{(1-A)\{Y-m_0(X)\}}{1-e(X)}.
\]
The independent-unit mean operation returns an estimate and variance
$N^{-2}\sum_i(U_i-\bar U)^2$, conditional on the identification and
nuisance-rate arguments in Chapter~\ref{r3:learning}. It rejects a declared
calendar-dependent sampling law. The score algebra is reusable; its sampling
covariance is not automatically portable to common shocks or repeated rows.
Independent path checks include a blocked confounder, an unblocked hidden
confounder, an opened collider, an action descendant and a cyclic graph.

A second adjustment query substitutes participation for wealth as the
outcome, producing first-stage scores $U_D$ alongside reduced-form scores
$U_Y$ on identical customers and splits. The ratio operation composes them:
\[
 \hat\theta=\bar U_Y/\bar U_D,
 \qquad \hat\phi_i=(U_{Y,i}-\hat\theta U_{D,i})/\bar U_D.
\]
It uses their joint covariance and residual-test inversion to retain weak
denominator topology. Instrument independence, exclusion, monotonicity and
positive complier mass are additional obligations before calling this ratio
a LATE. They do not follow from the graph's eligibility-to-wealth adjustment
trace. Bernoulli participation regressions are constrained to $[0,1]$;
that support projection is recorded separately from propensity clipping.

The public household-finance study and a separate synthetic bank contact/profit
study use the same query, graph, nuisance, score and mean operations with
different fields, adjustment dimension, populations and currencies. The
longitudinal query in Section~\ref{s5:history-censoring} has a different
identification expression and backwards nuisance recursion but reuses the
independent-unit mean/variance operation on customer trajectory scores.
The temporal query in Section~\ref{s5:discovery} reuses the principle of an
explicit target and result state, while returning conditional curve sets
rather than forcing graph ambiguity into a mean-score interface. These are
distinct reusable module compositions. They do not constitute a generic
compiler for every SCM or a complete application workbench.

\subsection{A substantive package adapter with a matched reference}
DoubleML 0.11.4 provides an installed IRM ATE operation
\citep{doubleml2022,runtime_doubleml_scope}. The adapter supplies the exact
same folds and nuisance predictions as the independent score calculation,
uses unnormalized inverse weights and propensity truncation at .01, and
checks the resulting estimate and standard error. Across the six public
learner/split runs and the second bank-contact example, coefficient
discrepancy is zero and standard-error discrepancy is at most
$4.55\times10^{-13}$. This is a meaningful test of target, weight and
covariance semantics under matched inputs. Comparing independently fitted
learners would conflate those semantics with nuisance variation; the public
study reports that second comparison separately.

The software choice is to reuse DoubleML for the qualified static
semiparametric operation and Tigramite for qualified temporal search. Build
the observation/clock contracts, identification traces, result-state
preservation and cross-library comparison layer. Keep the explicit references
for controlled likelihoods, dependence-sensitive moments and generated
observation plans, because these operations carry assumptions that no
generic estimator interface supplies. Defer unrestricted PAG identification,
general neural POMDP inference and distributed execution until their
scientific and operational contracts are qualified.

The installed-source records supplement the 44-project document survey;
they do not retroactively turn that whole survey into runtime qualification.
DoubleML's archived distribution declares BSD licensing. Tigramite's
installed metadata declares GPLv3 or later, so the packaging/distribution
decision must preserve those obligations; a mathematical adapter test is
not distribution clearance. The standalone design exposes these operations
through explicit requests and preserved artifacts. Sibling integrations
map existing domain fields to those requests, with adapters tested here;
accepted execution inside the sibling repositories remains separate work.


\section{One dynamic scientific brief, two product consumers}
\label{s6:composition}
The treasury brief asks for a one-month ten-basis-point price shift's
seven-month funding contribution for a fixed portfolio. Source records
contain candidate entity assignments and delayed balance revisions.
The information plan specifies an independent nuisance pilot, independently
audited pair matches, fixed calendar origins, the conditional price-innovation
restriction, and the common-calendar covariance rule. A causal result
requires all of these inputs; a regression formula alone cannot substitute
for the observation and sampling experiment.

The local construction first materializes the two candidate outcome
histories. Predecision features must be available by the decision date.
Realized outcomes may arrive later, but only revisions available by the
fixed retrospective extraction cutoff enter the evaluation view. These
are different clocks: a realized four-month outcome cannot be required to
exist before its treatment. The raw fixture includes later revisions whose
balances depend on the action. They are excluded by their arrival time;
business dates alone would leak them into the view. Missing horizons,
duplicated origin features and unavailable predecision information are
rejected. Conflicting source semantics remain unresolved before fitting.
The selected materializer accepts the account unit and received-day
availability only. Even a uniquely supplied business-day interpretation is
rejected as unsupported by this operation. Source-owner assertions are
simulated provenance inputs here; real authentication and semantic review
are not established by the fixture's trusted flag.

The same observation object feeds the nuisance fit, stratum-specific response
moments, inverse-linkage correction, calendar plus calibration covariance
and funding account derived in Sections~\ref{s6:funding}
and~\ref{s6:linked-histories}. Source and audit identities are preserved
alongside hashes of the actual pilot, evaluation and audit arrays. Target
units, horizon, financial weights, linkage law, sampling law and bandwidth
and nuisance-selection candidates also enter the analysis identity. The
supplied source assertions and retrospective extraction cutoff are bound
as well, even when a changed assertion leaves numeric arrays unchanged.
Changing bytes while retaining an old
source name invalidates reuse, just as changing the sampling law does.

The treasury consumer displays total funding contribution, its interval
and the response curve. An application builder consumes the same result
and divides by an externally fixed 480,000-account exposure. In the
720-origin reference, the estimate is 132,418 currency units, with interval
$[-149{,}326,414{,}162]$. The second consumer reports .276 per account and
interval $[-.311,.863]$, carrying exactly the same scientific identity.
This rescaling is valid because its denominator is fixed. Estimating or
selecting the account population from outcomes would need an additional
influence term and target definition.

Changing the externally supplied funding margin from .002 to .0015 is
supported: the response estimate is reused scientifically, while the
financial linear combination and its covariance are recomputed under a
new identity. An attempted balance-stock-to-cash-flow reinterpretation is
rejected because the price cost and balance retention account no longer
apply. An independent-branch sampling request is rejected by the calendar
backend. Submitting an old cached identity after any of these economic or
data changes is rejected before its interval can be displayed. These are
executed interface outcomes, not only proposed validation messages.

The reusable layer therefore owns observation and intervention semantics,
uncertainty composition and economic result transformations. Numerical
regression, array algebra and mature graph search remain delegated to
qualified library operations. The direct statsmodels comparison uses the
same residuals and covariance convention \citep{runtime_statsmodels_dynamic};
the independent linked-response reference checks the complete correction
and financial covariance. The prototype deliberately has no generic
compiler for arbitrary identification formulas. New score/likelihood
families need an explicit implementation with its own assumptions.

This is a local common library and two reference consumers. FinMarketing,
dataml and MacroFinance can adapt their observation and result interfaces
to these objects, but no sibling adoption is asserted. A bank workbench
would add data access, review interaction, entitlement and lifecycle
features around these scientifically qualified operations. The build choice
is to own the common scientific contract, reuse numerical libraries and
defer unsupported nonlinear/PAG procedures. The observed small-workload
cost does not justify a distributed platform merely to run these examples.

\subsection{A financial set is a different public result}
\label{s61:projection-product}
The same public operation accepts the projection procedure in
Section~\ref{s61:linkage-sets}. Its result carries a finite list of interval
components, confidence level and the separate audit/response coverage
conditions. A missing lower or upper endpoint explicitly means an infinite
endpoint; it never means a missing estimate to be filled by zero. The treasury
view shows the monetary set, while the application view divides each
component by the fixed positive account count. A disconnected set remains
disconnected after this transformation. There is no invented point estimate
or symmetric standard error when the information only supports a set.

The regular and projection procedures are selected explicitly, and the
procedure, confidence level, fixed error allocation and sum/difference
coordinates enter the scientific identity. A caller cannot choose the
narrower interval after inspecting both and inherit either procedure's
coverage claim. Displaying them as sensitivity comparisons is different
from a prespecified inferential choice. The candidate-coordinate projection
is retained as a mathematical challenger, not an automatic width selector.

The input domain is part of the calculation. The public operation checks
finite observations and economic quantities, complete response dimensions,
the declared driver/stratum basis, full pilot-training rank, both evaluation
strata, residual price variation, binary calibration records and an integer
bandwidth below the calendar length. Conflicting copies of the same outcome
revision and nonfinite extraction clocks are rejected. Regular inverse-map
inference also needs audit records in both strata. Projection can instead
use $[0,1]$ for an unaudited stratum, often producing an uninformative set.
An unsupported variance or nonfinite financial result cannot carry a
successful computed status. These checks verify numerical and structural
compatibility; they do not establish the source owner's truth, pilot
independence, the causal innovation restriction or audit representativeness.

An executed adversarial case removes all audits from one stratum. The regular
operation rejects the missing calibration, while projection returns the
appropriate set under its larger probability region. Other retained cases
change the inference procedure and confidence level, remove an evaluation
stratum, introduce nonbinary audit labels, break the pilot basis, and supply
conflicting revisions. Their failures identify the particular information
or procedure that must change before a result can be requested again.

\section{Three experiments, shared meaning and deliberate refusal}
\label{c7:journeys}
The preceding constructions support reuse because they expose compatible
scientific meanings, not because every problem is converted into the same
regression. A request declares its population, intervention normalization,
unit, information clock, sampling law, inference method and fixed monetary
scale. The persistent-funding request invokes the independent-pair audited
moment. The servicing request invokes the identified nonlinear model and a
policy fixed independently of evaluation. The discovery request retains its
latent class or its separately justified encouragement restriction.

Each result distinguishes a numeric calculation from its uncertainty claim.
The inference description contains the method, nominal level, finite-sample
or asymptotic interpretation, assumptions, numerical status and adverse
empirical evidence. Both a treasury account and a customer-facing application
retain that description. The latter may divide a set by a fixed external
account count, preserving disconnected components and unbounded endpoints.
An estimated exposure denominator instead requires a joint influence argument.

This requirement repairs a practical defect in the earlier regular funding
path: its public consumers dropped the scope field even though the underlying
calculation carried it. The revised request makes the inference method explicit
before fitting. The retained regular method discloses its adverse calibration;
neither it nor the projection is selected after comparing favorable widths.
Legacy callers preserve their original procedure while receiving the fuller
description. The change preserves the original numeric results and changes
what their recipients are told.

The local examples execute two eligible requests within each family and send
each result to both consumers. The second funding request changes horizon and
sample size; the second servicing request changes the fixed observable policy;
the second discovery request supplies valid encouragement information. These
are new scientific inputs and receive new identities. The identity binds the
request, data and implementation, but it does not certify the truth of a
supplied assumption.

Incompatible requests fail at meaningful boundaries. A pair-level persistent
map cannot borrow a calendar HAC covariance. A current audited hidden state
cannot become an input to the deployed policy. An account-date cannot replace
the independent customer unit. A selected graph cannot inherit confidence from
an outer structural union without carrying the union's information conditions.
Wrong clocks, empty populations and unsupported inference combinations are
rejected. These refusals express missing scientific arguments and information,
not a generic software inconvenience.

Direct expert composition can call the same qualified kernels and satisfy the
same unit, provenance, uncertainty and reporting obligations. The shared product
is valuable only if it helps preserve those obligations across repeated tasks.
The local matched computation compares identical public output bytes; it does
not measure analyst savings or establish acceptance by FinMarketing,
BayesFilter, MacroFinance or dataml. Their inspected interfaces remain concrete
integration proposals until their owners qualify and adopt the mappings.

